% Options for packages loaded elsewhere
\PassOptionsToPackage{unicode}{hyperref}
\PassOptionsToPackage{hyphens}{url}
\documentclass[
]{article}
\usepackage{xcolor}
\usepackage[margin=1in]{geometry}
\usepackage{amsmath,amssymb}
\setcounter{secnumdepth}{-\maxdimen} % remove section numbering
\usepackage{iftex}
\ifPDFTeX
  \usepackage[T1]{fontenc}
  \usepackage[utf8]{inputenc}
  \usepackage{textcomp} % provide euro and other symbols
\else % if luatex or xetex
  \usepackage{unicode-math} % this also loads fontspec
  \defaultfontfeatures{Scale=MatchLowercase}
  \defaultfontfeatures[\rmfamily]{Ligatures=TeX,Scale=1}
\fi
\usepackage{lmodern}
\ifPDFTeX\else
  % xetex/luatex font selection
\fi
% Use upquote if available, for straight quotes in verbatim environments
\IfFileExists{upquote.sty}{\usepackage{upquote}}{}
\IfFileExists{microtype.sty}{% use microtype if available
  \usepackage[]{microtype}
  \UseMicrotypeSet[protrusion]{basicmath} % disable protrusion for tt fonts
}{}
\makeatletter
\@ifundefined{KOMAClassName}{% if non-KOMA class
  \IfFileExists{parskip.sty}{%
    \usepackage{parskip}
  }{% else
    \setlength{\parindent}{0pt}
    \setlength{\parskip}{6pt plus 2pt minus 1pt}}
}{% if KOMA class
  \KOMAoptions{parskip=half}}
\makeatother
\usepackage{color}
\usepackage{fancyvrb}
\newcommand{\VerbBar}{|}
\newcommand{\VERB}{\Verb[commandchars=\\\{\}]}
\DefineVerbatimEnvironment{Highlighting}{Verbatim}{commandchars=\\\{\}}
% Add ',fontsize=\small' for more characters per line
\usepackage{framed}
\definecolor{shadecolor}{RGB}{248,248,248}
\newenvironment{Shaded}{\begin{snugshade}}{\end{snugshade}}
\newcommand{\AlertTok}[1]{\textcolor[rgb]{0.94,0.16,0.16}{#1}}
\newcommand{\AnnotationTok}[1]{\textcolor[rgb]{0.56,0.35,0.01}{\textbf{\textit{#1}}}}
\newcommand{\AttributeTok}[1]{\textcolor[rgb]{0.13,0.29,0.53}{#1}}
\newcommand{\BaseNTok}[1]{\textcolor[rgb]{0.00,0.00,0.81}{#1}}
\newcommand{\BuiltInTok}[1]{#1}
\newcommand{\CharTok}[1]{\textcolor[rgb]{0.31,0.60,0.02}{#1}}
\newcommand{\CommentTok}[1]{\textcolor[rgb]{0.56,0.35,0.01}{\textit{#1}}}
\newcommand{\CommentVarTok}[1]{\textcolor[rgb]{0.56,0.35,0.01}{\textbf{\textit{#1}}}}
\newcommand{\ConstantTok}[1]{\textcolor[rgb]{0.56,0.35,0.01}{#1}}
\newcommand{\ControlFlowTok}[1]{\textcolor[rgb]{0.13,0.29,0.53}{\textbf{#1}}}
\newcommand{\DataTypeTok}[1]{\textcolor[rgb]{0.13,0.29,0.53}{#1}}
\newcommand{\DecValTok}[1]{\textcolor[rgb]{0.00,0.00,0.81}{#1}}
\newcommand{\DocumentationTok}[1]{\textcolor[rgb]{0.56,0.35,0.01}{\textbf{\textit{#1}}}}
\newcommand{\ErrorTok}[1]{\textcolor[rgb]{0.64,0.00,0.00}{\textbf{#1}}}
\newcommand{\ExtensionTok}[1]{#1}
\newcommand{\FloatTok}[1]{\textcolor[rgb]{0.00,0.00,0.81}{#1}}
\newcommand{\FunctionTok}[1]{\textcolor[rgb]{0.13,0.29,0.53}{\textbf{#1}}}
\newcommand{\ImportTok}[1]{#1}
\newcommand{\InformationTok}[1]{\textcolor[rgb]{0.56,0.35,0.01}{\textbf{\textit{#1}}}}
\newcommand{\KeywordTok}[1]{\textcolor[rgb]{0.13,0.29,0.53}{\textbf{#1}}}
\newcommand{\NormalTok}[1]{#1}
\newcommand{\OperatorTok}[1]{\textcolor[rgb]{0.81,0.36,0.00}{\textbf{#1}}}
\newcommand{\OtherTok}[1]{\textcolor[rgb]{0.56,0.35,0.01}{#1}}
\newcommand{\PreprocessorTok}[1]{\textcolor[rgb]{0.56,0.35,0.01}{\textit{#1}}}
\newcommand{\RegionMarkerTok}[1]{#1}
\newcommand{\SpecialCharTok}[1]{\textcolor[rgb]{0.81,0.36,0.00}{\textbf{#1}}}
\newcommand{\SpecialStringTok}[1]{\textcolor[rgb]{0.31,0.60,0.02}{#1}}
\newcommand{\StringTok}[1]{\textcolor[rgb]{0.31,0.60,0.02}{#1}}
\newcommand{\VariableTok}[1]{\textcolor[rgb]{0.00,0.00,0.00}{#1}}
\newcommand{\VerbatimStringTok}[1]{\textcolor[rgb]{0.31,0.60,0.02}{#1}}
\newcommand{\WarningTok}[1]{\textcolor[rgb]{0.56,0.35,0.01}{\textbf{\textit{#1}}}}
\usepackage{longtable,booktabs,array}
\usepackage{calc} % for calculating minipage widths
% Correct order of tables after \paragraph or \subparagraph
\usepackage{etoolbox}
\makeatletter
\patchcmd\longtable{\par}{\if@noskipsec\mbox{}\fi\par}{}{}
\makeatother
% Allow footnotes in longtable head/foot
\IfFileExists{footnotehyper.sty}{\usepackage{footnotehyper}}{\usepackage{footnote}}
\makesavenoteenv{longtable}
\usepackage{graphicx}
\makeatletter
\newsavebox\pandoc@box
\newcommand*\pandocbounded[1]{% scales image to fit in text height/width
  \sbox\pandoc@box{#1}%
  \Gscale@div\@tempa{\textheight}{\dimexpr\ht\pandoc@box+\dp\pandoc@box\relax}%
  \Gscale@div\@tempb{\linewidth}{\wd\pandoc@box}%
  \ifdim\@tempb\p@<\@tempa\p@\let\@tempa\@tempb\fi% select the smaller of both
  \ifdim\@tempa\p@<\p@\scalebox{\@tempa}{\usebox\pandoc@box}%
  \else\usebox{\pandoc@box}%
  \fi%
}
% Set default figure placement to htbp
\def\fps@figure{htbp}
\makeatother
\setlength{\emergencystretch}{3em} % prevent overfull lines
\providecommand{\tightlist}{%
  \setlength{\itemsep}{0pt}\setlength{\parskip}{0pt}}
\usepackage{bookmark}
\IfFileExists{xurl.sty}{\usepackage{xurl}}{} % add URL line breaks if available
\urlstyle{same}
\hypersetup{
  pdftitle={Evolutionary rescue under habitat fragmentation},
  pdfauthor={Jan Leuenberger \& Safeeya Ramis},
  hidelinks,
  pdfcreator={LaTeX via pandoc}}

\title{Evolutionary rescue under habitat fragmentation}
\author{Jan Leuenberger \& Safeeya Ramis}
\date{2026-10-28}

\begin{document}
\maketitle

\section*{Abstract}\label{abstract}
\addcontentsline{toc}{section}{Abstract}

\section{Introduction}\label{introduction}

The Anthropocene is characterized by a wide variety of changes and
threats that affect nearly all wildlife populations. These human-induced
pressures are having an increasingly negative impact on the stability
and growth of populations and are leading to a continuous decline in the
individual numbers within the populations. In times like these, the idea
of \textbf{evolutionary rescue} (ER) is becoming increasingly important.
\emph{``ER occurs when genetic adaptation allows a population to recover
from demographic effects initiated by environmental change that would
otherwise cause extirpation''} (Gonzalez et al., 2013) and therefore
leads back to an increase in the number of individuals within a
population.

Because ER is often viewed as a solution to environmental change and
pressure, we thought about what influence a major environmental threat
such as \textbf{habitat fragmentation}; \emph{``the process by which
habitats are transformed into smaller patches isolated from each other
{[}\ldots{]}''} (Cheptou et al., 2017) might have on an ER process that
has already begun.

\subsection{Research Questions}\label{research-questions}

Based on this reflection, the following questions emerged, which we
intended to explore in more detail as part of a three-week theoretical
practical course at the University of Bern:

\begin{itemize}
\item
  Is evolutionary rescue possible in fragmented habitats if the habitats
  are separated at a certain time point after the ER processes have
  already begun?

  \begin{itemize}
  \item
    How does the timing of habitat fragmentation affect evolutionary
    rescue?
  \item
    How does the size of the habitats affect evolutionary rescue?
  \item
    What impact does the distribution of individuals across habitats
    have on evolutionary rescue?
  \item
    What would happen if migration between the isolated fragmented
    habitats were possible?
  \end{itemize}
\end{itemize}

It is important to analyze these questions in more detail because
habitat fragmentation represents an increasing threat to populations,
which can occur at any time. Furthermore, from a conservation
perspective, it is important to understand the impact of splitting time,
habitat size, and individual distribution on evolutionary rescue in
order to ensure the survival of the populations.

This report will document a stochastic simulation of
\textbf{evolutionary rescue} in discrete time, based on the model of Orr
\& Unckless introduced in the practical, which was extended with
different parameters to answer the questions listed above. Therefore, an
explanation of the model is provided in the Methods section, followed by
a presentation of the results obtained and a subsequent discussion.

\section{Methodology}\label{methodology}

\subsection{Base Model}\label{base-model}

The base model of Orr \& Unckless, which was provided from the banklab
recreating a biological scenario on which we built on. In this scenario
a haploid population consists of wild-type individuals (\textbf{a}) and,
possibly, mutant individuals (\textbf{A}). A sudden environmental change
has a negative efefect on the wild type and makes it unfit, which is why
it starts to \textbf{decline}. The mutant allele is \textbf{beneficial}
and is able to grow under the new environmental conditions, but it is
initially rare or absent. If it's absent, it has to arise first by
mutation from the wild type. The population is then \textbf{rescued} if
the mutant spreads before the wild type goes extinct.

\subsection{Fragmentation and
migration}\label{fragmentation-and-migration}

To assess the impact of \textbf{habitat fragmentation} on ER, we
extended the model by adding a additional biological scenario, where the
habitat, in which the population is located (called habitat A), gets
separated after a defined number of generations into \textbf{two
isolated habitats} (called habitat B and C) of defined sizes. Based on
that, the following equation must be considered with regard to the
habitat sizes:\((A_size = B_size + C_size)\). Each habitat has also its
own carrying capacity (\textbf{K}) related to its size:
\((K_A = K_B + K_C)\). At a \textbf{separation time} individuals of the
population A are distributed then into the two new habitats B and C and
both will then evolve independently. Since our aim was also to
investigate what happens when \textbf{migration} is possible, we allowed
individuals to move between habitats B and C after the splitting event.
In each generation, every individual (wild type and mutant separately)
migrates to the other habitat with probability \textbf{migrate}.

\subsection{Parameters and variables}\label{parameters-and-variables}

Overall, we have expanded the base model therefore with five parameters;
\texttt{split\_time,\ habitat\_prop,\ split\_ratio,\ K} and
\texttt{migration}. All the other parameters and variables were already
set in the based model. The table below provides brief explanations of
the respective parameters and variables.

\textbf{Taken from the base model:}

\begin{longtable}[]{@{}
  >{\raggedright\arraybackslash}p{(\linewidth - 2\tabcolsep) * \real{0.5000}}
  >{\raggedright\arraybackslash}p{(\linewidth - 2\tabcolsep) * \real{0.5000}}@{}}
\toprule\noalign{}
\begin{minipage}[b]{\linewidth}\raggedright
Name
\end{minipage} & \begin{minipage}[b]{\linewidth}\raggedright
Meaning
\end{minipage} \\
\midrule\noalign{}
\endhead
\bottomrule\noalign{}
\endlastfoot
\texttt{N\_a}, \texttt{N\_A} & current size of the wild-type / mutant
population (change over time) \\
\texttt{N\_init\_a}, \texttt{N\_init\_A} & initial sizes of wild type
(positive integer) and mutant (non-negative integer) \\
\texttt{decay\_rate} & detrimental effect of the environment on the wild
type (\textgreater{} 0 so that extinction is possible) \\
\texttt{sel\_coeff} & selective advantage of the mutant (should be
\textgreater{} \texttt{decay\_rate} to allow rescue) \\
\texttt{mut\_rate} & mutation rate from wild type to mutant
(\textgreater{} 0 and small) \\
\texttt{t\_max} & maximum number of generations simulated (a setting,
not a model parameter) \\
\end{longtable}

\textbf{Added to expand the model:}

\begin{longtable}[]{@{}
  >{\raggedright\arraybackslash}p{(\linewidth - 2\tabcolsep) * \real{0.5000}}
  >{\raggedright\arraybackslash}p{(\linewidth - 2\tabcolsep) * \real{0.5000}}@{}}
\toprule\noalign{}
\begin{minipage}[b]{\linewidth}\raggedright
Name
\end{minipage} & \begin{minipage}[b]{\linewidth}\raggedright
Meaning
\end{minipage} \\
\midrule\noalign{}
\endhead
\bottomrule\noalign{}
\endlastfoot
\texttt{split\_time} & generation at which the fragmentation happens
(integer \textgreater{} 0) \\
\texttt{habitat\_prop} & defining size of the new habitat B and
therefore also it's carrying capacity (between 0 and 1); habitat C gets
\texttt{1\ -\ habitat\_prop} \\
\texttt{split\_ratio} & probability with which each individual is
assigned to habitat B at the split (between 0 and 1): habitat C gets
\texttt{1\ -\ split\_ratio} \\
\texttt{K} & carrying capacity of the original (unsplit) habitat; limits
growth and is divided among B and C in the same proportions as the
habitat size \\
\texttt{migrate} & migration rate: (between 0 and 1) with which each
individual moves to the other habitat in every generation after the
split. \\
\end{longtable}

\subsubsection{Example}\label{example}

To ensure that there is no confusion and that the parameters and
variables are clear, here is a brief example:

\begin{longtable}[]{@{}ll@{}}
\toprule\noalign{}
Parameters & values \\
\midrule\noalign{}
\endhead
\bottomrule\noalign{}
\endlastfoot
\texttt{N\_a}, \texttt{N\_A} & 200, 100 \\
\texttt{split\_time} & 30 \\
\texttt{habitat\_prop} & 0.8 \\
\texttt{split\_ratio} & 0.7 \\
\texttt{K} & 1000 \\
\texttt{migrate} & 0.01 \\
\end{longtable}

The carrying capacity of habitat A is set at 1000 individuals. Based on
the \texttt{habitat\_prop} we know that habitat B is able to carry 800
individuals \((K * habitat_{prop})\) and habitat C is able to carry 200
individuals \((K*(1-habitat_{prop}))\). From the \texttt{split\_ratio}
we know that 140 wild type individuals and 70 mutant individuals are
located within the habitat B \((N_{a} * split_{time})\) \&
\((N_{A} * split_{time})\). In habitat C we will then be able to find 60
wild type individuals and 30 mutant individuals
\((N_{a} * (1- split_{time})\) \& \((N_{A} * (1- split_{time}))\). With
\texttt{migrate} = 0.01, each individual has a 1\% chance per generation
of moving to the other habitat.

\subsection{Model/Simulation}\label{modelsimulation}

\subsubsection{Reproducibility}\label{reproducibility}

To ensure the reproducibility of our results, it is important to set a
seed first, in our case 57. In the same chunk we installed the libraries
needed to run our simulation and analysis.

\subsubsection{Overview of the code}\label{overview-of-the-code}

\begin{enumerate}
\def\labelenumi{\arabic{enumi}.}
\tightlist
\item
  \texttt{simulate\_one\_gen()}: simulates \textbf{one generation}
  (reproduction with logistic growth + mutation) for one habitat.
\item
  \texttt{simulate\_pop()}: runs a \textbf{full simulation}: one habitat
  until \texttt{split\_time}, then two habitats (B and C) afterwards.
\item
  Plot functions (\texttt{plot\_split()}, \texttt{plot\_habitat()},
  \texttt{plot\_habitats()}): visualise a single simulation run.
\item
  \texttt{run\_scenarios()} and \texttt{to\_long()}: repeat the
  simulation many times for different parameter values and summarise the
  \textbf{rescue probability}.
\item
  Six analyses: split time varies, habitat proportion varies and
  migration rates
\end{enumerate}

\subsubsection{Simulate one generation}\label{simulate-one-generation}

Our actual simulation begins in the second chunk, although the change in
population is first calculated only within a single generation. To do
this, we use the function \texttt{simulate\_one\_gen()}, which simulates
\textbf{one generation} based on the following steps:

\begin{enumerate}
\def\labelenumi{\arabic{enumi}.}
\item
  \textbf{Expected offspring (logistic growth)}

  Without competition, the per-capita growth rates of the individual
  are:

  \begin{itemize}
  \item
    wild type: \(\rho_a = -r\) (negative, so the type declines), with
    \(r\) = \texttt{decay\_rate}
  \item
    mutant: \(\rho_A = s - r\) (positive if \(s > r\)), with \(s\) =
    \texttt{sel\_coeff}
  \end{itemize}

  In addition a carrying capacity is added by a
  \textbf{density-dependent term}, means that the growth factor of each
  type is reduced by \(\rho_A \cdot N_{tot}/K\). The expected number of
  offspring is therefore

  \[\bar N_a = N_a\,\max\!\left(0,\;1+\rho_a-\rho_A \tfrac{N_{tot}}{K}\right), \qquad \bar N_A = N_A\,\max\!\left(0,\;1+\rho_A-\rho_A \tfrac{N_{tot}}{K}\right)\]

  When \(N_{tot}\) approaches \(K\), the mutant's growth factor
  approaches 1 (population stays constant) and the wild type keeps
  declining. When \(N_{tot} > K\) the mutant would shrink back towards
  \(K\). The \texttt{max(0,\ ...)} prevents negative expected values,
  which \texttt{rpois()} cannot handle.
\item
  \textbf{Demographic stochasticity}

  The actual number of offspring is drawn from a \textbf{Poisson
  distribution} whose mean is the expected value above. This is what
  makes every run different.
\item
  \textbf{Mutations}

  Of the wild-type offspring (a), a Poisson-distributed number (mean
  \texttt{offsp\_a\ *\ mut\_rate}) mutates to mutant individuals (A).
  These are subtracted from the wild type and added to the mutants.
\end{enumerate}

The function then returns a vector \texttt{c(N\_a\_new,\ N\_A\_new)},
containing the current population size based on these calculations.

Since the base model describes exponential growth and this needed to be
converted to logistic growth in this model, we used AI to some extent to
adapt the code for this purpose.

\begin{Shaded}
\begin{Highlighting}[]
\NormalTok{simulate\_one\_gen }\OtherTok{\textless{}{-}} \ControlFlowTok{function}\NormalTok{(N\_a, N\_A, decay\_rate, sel\_coeff, mut\_rate, K) \{}
  
  \CommentTok{\# total population size and per{-}capita growth rates without competition}
\NormalTok{  N\_tot }\OtherTok{\textless{}{-}}\NormalTok{ N\_a }\SpecialCharTok{+}\NormalTok{ N\_A}
\NormalTok{  rho\_a }\OtherTok{\textless{}{-}} \SpecialCharTok{{-}}\NormalTok{decay\_rate}
\NormalTok{  rho\_A }\OtherTok{\textless{}{-}}\NormalTok{ sel\_coeff }\SpecialCharTok{{-}}\NormalTok{ decay\_rate}
  \CommentTok{\# logistic growth: crowding reduces the growth of both types equally (set to 0 if negative)}
\NormalTok{  mean\_a }\OtherTok{\textless{}{-}}\NormalTok{ N\_a }\SpecialCharTok{*} \FunctionTok{max}\NormalTok{(}\DecValTok{0}\NormalTok{, }\DecValTok{1} \SpecialCharTok{+}\NormalTok{ rho\_a }\SpecialCharTok{{-}}\NormalTok{ rho\_A }\SpecialCharTok{*}\NormalTok{ N\_tot }\SpecialCharTok{/}\NormalTok{ K)}
\NormalTok{  mean\_A }\OtherTok{\textless{}{-}}\NormalTok{ N\_A }\SpecialCharTok{*} \FunctionTok{max}\NormalTok{(}\DecValTok{0}\NormalTok{, }\DecValTok{1} \SpecialCharTok{+}\NormalTok{ rho\_A }\SpecialCharTok{{-}}\NormalTok{ rho\_A }\SpecialCharTok{*}\NormalTok{ N\_tot }\SpecialCharTok{/}\NormalTok{ K)}
  \CommentTok{\# draw offspring according to Poisson distribution}
\NormalTok{  offsp\_a }\OtherTok{\textless{}{-}} \FunctionTok{rpois}\NormalTok{(}\DecValTok{1}\NormalTok{, mean\_a)}
\NormalTok{  offsp\_A }\OtherTok{\textless{}{-}} \FunctionTok{rpois}\NormalTok{(}\DecValTok{1}\NormalTok{, mean\_A)}
  \CommentTok{\# draw new mutants according to Poisson distribution}
\NormalTok{  mut\_a\_to\_A }\OtherTok{\textless{}{-}} \FunctionTok{rpois}\NormalTok{(}\DecValTok{1}\NormalTok{, offsp\_a }\SpecialCharTok{*}\NormalTok{ mut\_rate)}
  \CommentTok{\# determine new population sizes of wild type and mutant}
\NormalTok{  N\_a\_new }\OtherTok{\textless{}{-}}\NormalTok{ offsp\_a }\SpecialCharTok{{-}}\NormalTok{ mut\_a\_to\_A}
\NormalTok{  N\_A\_new }\OtherTok{\textless{}{-}}\NormalTok{  offsp\_A }\SpecialCharTok{+}\NormalTok{ mut\_a\_to\_A}
  
  \FunctionTok{return}\NormalTok{(}\FunctionTok{c}\NormalTok{(N\_a\_new, N\_A\_new))}
\NormalTok{\}}
\CommentTok{\# Test the function}
\CommentTok{\#print(simulate\_one\_gen(100,0,0.1,0.2,0.001,1000))}
\end{Highlighting}
\end{Shaded}

\subsubsection{Simulating the whole
populations}\label{simulating-the-whole-populations}

The function \texttt{simulate\_one\_gen()} describes a single
generation, whereas evolutionary rescue unfolds over many. We therefore
iterate this step in \texttt{simulate\_pop()}, which integrates habitat
fragmentation (the population is divided into two isolated habitats at a
chosen generation) and migration (individuals may move between these
habitats afterwards). The function records the numbers of wild-type and
mutant individuals in each habitat, so we can determine whether the
mutant spreads and rescues the population, and how the timing of the
split, the relative habitat sizes and the migration rate influence this
outcome.

The function \texttt{simulate\_pop()} runs the full simulation over
\texttt{t\_max} generations and has three phases:

\begin{enumerate}
\def\labelenumi{\arabic{enumi}.}
\item
  \textbf{generations (i) \textless{} \texttt{split\_time\ +\ 1}} In
  this first phase, the population lives in one habitat with carrying
  capacity K, and each generation is simulated with
  \texttt{simulate\_one\_gen()}. This phase stops early if the
  population goes extinct or grows to at least twice its initial size.
\item
  \textbf{generation (i) == \texttt{split\_time\ +\ 1}} In the second
  phase, at generation \texttt{split\_time\ +\ 1}, the habitat is
  fragmented into two isolated habitats, B and C. The habitat sizes are
  calculated based on the \texttt{habitat\_prop}, as explained above.
  The existing individuals are distributed as well based on the
  \texttt{split\_ratio} using \texttt{rbinom()}.
\item
  \textbf{generation (i) \textgreater{} \texttt{split\_time\ +\ 1}} From
  then on both habitats are simulated independently, each with its own
  carrying capacity. After each generation, migration can occur, where
  individual moves to the other habitat with probability migrate, drawn
  separately for wild types and mutants. Migration is symmetric in its
  rate but not in its effect. The number of migrants is proportional to
  the size of the source habitat and migration is possible in both
  directions.
\end{enumerate}

It's important to note that \texttt{spit\_time} must be equal to 1,
because in the for loop, \texttt{i} isn't counted the same way as our
number of generations, but is shifted by one. Furthermore, the three
phases mentioned above are not represented in the same order in the
code. However, this does not affect the calculations, since the if
statements check the number of past generations.

At the end the function returns a list of three data frames: main
(habitat A), B and C, containing the population sizes for each habitat.
If the population dies out or reaches twice its starting size before the
split, B and C are empty.

\begin{Shaded}
\begin{Highlighting}[]
\NormalTok{simulate\_pop }\OtherTok{\textless{}{-}} \ControlFlowTok{function}\NormalTok{(N\_init\_a, N\_init\_A, decay\_rate, sel\_coeff, mut\_rate, t\_max, K, habitat\_prop, split\_ratio, split\_time,migrate) \{}
  \CommentTok{\# Create the vector in which to save the results}
\NormalTok{  pop\_vector }\OtherTok{\textless{}{-}} \FunctionTok{c}\NormalTok{(N\_init\_a, N\_init\_A)}
  \CommentTok{\# initiate the variables}
\NormalTok{  pop\_new }\OtherTok{\textless{}{-}} \FunctionTok{c}\NormalTok{(N\_init\_a, N\_init\_A)}
  \CommentTok{\# Create the vectors for the new habitats}
\NormalTok{  pop\_vector\_B }\OtherTok{\textless{}{-}} \ConstantTok{NULL}
\NormalTok{  pop\_vector\_C }\OtherTok{\textless{}{-}} \ConstantTok{NULL}
  \CommentTok{\# run the simulation until generation t\_max}
  \ControlFlowTok{for}\NormalTok{ (i }\ControlFlowTok{in} \DecValTok{1}\SpecialCharTok{:}\NormalTok{t\_max) \{}
    \CommentTok{\# Time when the fragmentation occur}
    \ControlFlowTok{if}\NormalTok{ (i }\SpecialCharTok{==}\NormalTok{ split\_time }\SpecialCharTok{+} \DecValTok{1}\NormalTok{) \{}
      \CommentTok{\# Create the new habitats (creating new K)}
\NormalTok{      K\_B }\OtherTok{\textless{}{-}}\NormalTok{ K }\SpecialCharTok{*}\NormalTok{ habitat\_prop}
\NormalTok{      K\_C }\OtherTok{\textless{}{-}}\NormalTok{ K }\SpecialCharTok{*}\NormalTok{ (}\DecValTok{1} \SpecialCharTok{{-}}\NormalTok{ habitat\_prop)}
      \CommentTok{\#create the ratio of the splitting}
     
\NormalTok{      split\_ratio\_B }\OtherTok{\textless{}{-}}\NormalTok{ split\_ratio}
\NormalTok{      split\_ratio\_C }\OtherTok{\textless{}{-}}\NormalTok{ (}\DecValTok{1} \SpecialCharTok{{-}}\NormalTok{ split\_ratio\_B)}
     
      \CommentTok{\# Distribute individuals  randomly to the habitat B}
\NormalTok{      n\_B\_a }\OtherTok{\textless{}{-}} \FunctionTok{rbinom}\NormalTok{(}\DecValTok{1}\NormalTok{, pop\_new[}\DecValTok{1}\NormalTok{], split\_ratio\_B) }
\NormalTok{      n\_B\_A }\OtherTok{\textless{}{-}} \FunctionTok{rbinom}\NormalTok{(}\DecValTok{1}\NormalTok{, pop\_new[}\DecValTok{2}\NormalTok{], split\_ratio\_B)   }
     
      \CommentTok{\# Distribute individuals randomly to the habitat C}
\NormalTok{      n\_C\_a }\OtherTok{\textless{}{-}}\NormalTok{ pop\_new[}\DecValTok{1}\NormalTok{] }\SpecialCharTok{{-}}\NormalTok{ n\_B\_a}
\NormalTok{      n\_C\_A }\OtherTok{\textless{}{-}}\NormalTok{ pop\_new[}\DecValTok{2}\NormalTok{] }\SpecialCharTok{{-}}\NormalTok{ n\_B\_A}
     
      \CommentTok{\# Calculate the new populations for B and C }
\NormalTok{      pop\_new\_B }\OtherTok{\textless{}{-}} \FunctionTok{simulate\_one\_gen}\NormalTok{(n\_B\_a, n\_B\_A, decay\_rate, sel\_coeff, mut\_rate, K\_B)}
\NormalTok{      pop\_new\_C }\OtherTok{\textless{}{-}} \FunctionTok{simulate\_one\_gen}\NormalTok{(n\_C\_a, n\_C\_A, decay\_rate, sel\_coeff, mut\_rate, K\_C)}
      
      \CommentTok{\# add the new population sizes to the output vector}
\NormalTok{      pop\_vector\_B }\OtherTok{\textless{}{-}} \FunctionTok{rbind}\NormalTok{(pop\_vector\_B,pop\_new\_B)}
\NormalTok{      pop\_vector\_C }\OtherTok{\textless{}{-}} \FunctionTok{rbind}\NormalTok{(pop\_vector\_C,pop\_new\_C)}
\NormalTok{    \}}
    
    \ControlFlowTok{if}\NormalTok{ (i }\SpecialCharTok{\textgreater{}}\NormalTok{ split\_time }\SpecialCharTok{+} \DecValTok{1}\NormalTok{) \{}
      
\NormalTok{      pop\_new\_B }\OtherTok{\textless{}{-}} \FunctionTok{simulate\_one\_gen}\NormalTok{(}
\NormalTok{        pop\_new\_B[}\DecValTok{1}\NormalTok{], pop\_new\_B[}\DecValTok{2}\NormalTok{],}
\NormalTok{        decay\_rate, sel\_coeff, mut\_rate,K\_B)}
      
\NormalTok{      pop\_new\_C }\OtherTok{\textless{}{-}} \FunctionTok{simulate\_one\_gen}\NormalTok{(}
\NormalTok{        pop\_new\_C[}\DecValTok{1}\NormalTok{], pop\_new\_C[}\DecValTok{2}\NormalTok{],}
\NormalTok{        decay\_rate, sel\_coeff, mut\_rate,K\_C)}
      
      \CommentTok{\# Migration }
\NormalTok{      migrateB\_vec }\OtherTok{\textless{}{-}} \FunctionTok{c}\NormalTok{(}\FunctionTok{rbinom}\NormalTok{(}\DecValTok{1}\NormalTok{, pop\_new\_B[}\DecValTok{1}\NormalTok{], migrate),}
                        \FunctionTok{rbinom}\NormalTok{(}\DecValTok{1}\NormalTok{, pop\_new\_B[}\DecValTok{2}\NormalTok{], migrate))}
\NormalTok{      migrateC\_vec }\OtherTok{\textless{}{-}} \FunctionTok{c}\NormalTok{(}\FunctionTok{rbinom}\NormalTok{(}\DecValTok{1}\NormalTok{, pop\_new\_C[}\DecValTok{1}\NormalTok{], migrate),}
                        \FunctionTok{rbinom}\NormalTok{(}\DecValTok{1}\NormalTok{, pop\_new\_C[}\DecValTok{2}\NormalTok{], migrate))}
      
      \CommentTok{\#adding the migrating individuals back to the Populations}
\NormalTok{      pop\_new\_B }\OtherTok{\textless{}{-}}\NormalTok{ pop\_new\_B }\SpecialCharTok{{-}}\NormalTok{ migrateB\_vec }\SpecialCharTok{+}\NormalTok{ migrateC\_vec}
\NormalTok{      pop\_new\_C }\OtherTok{\textless{}{-}}\NormalTok{ pop\_new\_C }\SpecialCharTok{{-}}\NormalTok{ migrateC\_vec }\SpecialCharTok{+}\NormalTok{ migrateB\_vec}

      \CommentTok{\# Save the results for each generation}
\NormalTok{      pop\_vector\_B }\OtherTok{\textless{}{-}} \FunctionTok{rbind}\NormalTok{(pop\_vector\_B, pop\_new\_B)}
\NormalTok{      pop\_vector\_C }\OtherTok{\textless{}{-}} \FunctionTok{rbind}\NormalTok{(pop\_vector\_C, pop\_new\_C)}
      
      \CommentTok{\# define the row and column names of the output vector}
      \FunctionTok{rownames}\NormalTok{(pop\_vector\_B) }\OtherTok{\textless{}{-}}\NormalTok{ (}\DecValTok{0}\SpecialCharTok{:}\NormalTok{t\_max)[}\DecValTok{1}\SpecialCharTok{:}\FunctionTok{length}\NormalTok{(pop\_vector\_B[,}\DecValTok{1}\NormalTok{])] }
      \FunctionTok{rownames}\NormalTok{(pop\_vector\_C) }\OtherTok{\textless{}{-}}\NormalTok{ (}\DecValTok{0}\SpecialCharTok{:}\NormalTok{t\_max)[}\DecValTok{1}\SpecialCharTok{:}\FunctionTok{length}\NormalTok{(pop\_vector\_C[,}\DecValTok{1}\NormalTok{])]}
      \FunctionTok{colnames}\NormalTok{(pop\_vector\_B) }\OtherTok{\textless{}{-}} \FunctionTok{c}\NormalTok{(}\StringTok{"a"}\NormalTok{,}\StringTok{"A"}\NormalTok{)}
      \FunctionTok{colnames}\NormalTok{(pop\_vector\_C) }\OtherTok{\textless{}{-}} \FunctionTok{c}\NormalTok{(}\StringTok{"a"}\NormalTok{,}\StringTok{"A"}\NormalTok{)}
      
      \ControlFlowTok{if}\NormalTok{ (}\FunctionTok{sum}\NormalTok{(pop\_new\_B) }\SpecialCharTok{+} \FunctionTok{sum}\NormalTok{(pop\_new\_C) }\SpecialCharTok{\textgreater{}=} \FloatTok{2.5}\SpecialCharTok{*}\NormalTok{(N\_init\_a }\SpecialCharTok{+}\NormalTok{ N\_init\_A) }\SpecialCharTok{|}
          \FunctionTok{sum}\NormalTok{(pop\_new\_B) }\SpecialCharTok{+} \FunctionTok{sum}\NormalTok{(pop\_new\_C) }\SpecialCharTok{==} \DecValTok{0}\NormalTok{) }\ControlFlowTok{break}
\NormalTok{    \}}
    
    \CommentTok{\# Continue simulating habitat A only until split\_time}
    \ControlFlowTok{if}\NormalTok{ (i }\SpecialCharTok{\textless{}}\NormalTok{ split\_time }\SpecialCharTok{+} \DecValTok{1}\NormalTok{) \{}
      
      \CommentTok{\# redefine the current population one generation later}
\NormalTok{      pop\_new }\OtherTok{\textless{}{-}} \FunctionTok{simulate\_one\_gen}\NormalTok{(pop\_new[}\DecValTok{1}\NormalTok{],pop\_new[}\DecValTok{2}\NormalTok{], decay\_rate, sel\_coeff, mut\_rate,K)}
      \CommentTok{\# add the new population sizes to the output vector}
\NormalTok{      pop\_vector }\OtherTok{\textless{}{-}} \FunctionTok{rbind}\NormalTok{(pop\_vector,pop\_new)}
      \CommentTok{\# condition to stop the simulation before t\_max: either the population exceeds 2 times the original size, or it goes extinct}
      \ControlFlowTok{if}\NormalTok{ (pop\_new[}\DecValTok{1}\NormalTok{]}\SpecialCharTok{+}\NormalTok{pop\_new[}\DecValTok{2}\NormalTok{]}\SpecialCharTok{\textgreater{}=}\DecValTok{2}\SpecialCharTok{*}\NormalTok{(N\_init\_a}\SpecialCharTok{+}\NormalTok{ N\_init\_A) }\SpecialCharTok{|}\NormalTok{ pop\_new[}\DecValTok{1}\NormalTok{]}\SpecialCharTok{+}\NormalTok{pop\_new[}\DecValTok{2}\NormalTok{]}\SpecialCharTok{==}\DecValTok{0}\NormalTok{) }\ControlFlowTok{break}
\NormalTok{    \}}
    
\NormalTok{  \}}
  \CommentTok{\# define the row and column names of the output vector}
  \FunctionTok{rownames}\NormalTok{(pop\_vector) }\OtherTok{\textless{}{-}}\NormalTok{ (}\DecValTok{0}\SpecialCharTok{:}\NormalTok{t\_max)[}\DecValTok{1}\SpecialCharTok{:}\FunctionTok{length}\NormalTok{(pop\_vector[,}\DecValTok{1}\NormalTok{])] }\CommentTok{\# note that the vector has to be cut if the simulation stopped early}
  \FunctionTok{colnames}\NormalTok{(pop\_vector) }\OtherTok{\textless{}{-}} \FunctionTok{c}\NormalTok{(}\StringTok{"a"}\NormalTok{,}\StringTok{"A"}\NormalTok{)}
  \CommentTok{\# Data frames}
\NormalTok{  pop\_vector\_df }\OtherTok{\textless{}{-}} \FunctionTok{as.data.frame}\NormalTok{(pop\_vector)}
\NormalTok{  pop\_vector\_B\_df }\OtherTok{\textless{}{-}} \FunctionTok{as.data.frame}\NormalTok{(pop\_vector\_B)}
\NormalTok{  pop\_vector\_C\_df }\OtherTok{\textless{}{-}} \FunctionTok{as.data.frame}\NormalTok{(pop\_vector\_C)}
  \CommentTok{\# return all three population vectors separately}
  \FunctionTok{return}\NormalTok{(}\FunctionTok{list}\NormalTok{(}\AttributeTok{main =}\NormalTok{ pop\_vector\_df, }\AttributeTok{B =}\NormalTok{ pop\_vector\_B\_df, }\AttributeTok{C =}\NormalTok{ pop\_vector\_C\_df))}
\NormalTok{\}}
\end{Highlighting}
\end{Shaded}

\subsubsection{Plotting functions for the
model}\label{plotting-functions-for-the-model}

To inspect what happens in a single simulation, we define three plot
functions. plot\_split() shows the whole run: wild type and mutant
before the split, and the total population of habitats B and C
afterwards, together with the carrying capacities (K, K\_B, K\_C) and
the split time. plot\_habitat() draws one habitat in detail, with the
wild type, the mutant and the total population, and plot\_habitats()
stacks the panels of B and C so that they can be compared directly.
These plots serve to check that the simulation behaves as expected, and
they were used extensively for debugging in the early stages of the
project. We used base R graphics here, and the plots received very
positive feedback. Opinions on whether to use base R or ggplot2 for
plotting are divided, even within the \textbf{THEE} Lab. In the later
plotting functions we use ggplot2, both as an exercise to learn the
package and to show that both approaches can produce informative
figures.

\begin{Shaded}
\begin{Highlighting}[]
\CommentTok{\# ============================================================}
\CommentTok{\# Usage}
\CommentTok{\# (simulate\_one\_gen and simulate\_pop must be defined first)}
\CommentTok{\# ============================================================}
\NormalTok{max\_gen      }\OtherTok{\textless{}{-}} \DecValTok{1000}
\NormalTok{K            }\OtherTok{\textless{}{-}} \DecValTok{1000}
\NormalTok{habitat\_prop }\OtherTok{\textless{}{-}} \FloatTok{0.8}
\NormalTok{split\_time   }\OtherTok{\textless{}{-}} \DecValTok{50}
\NormalTok{split\_ratio  }\OtherTok{\textless{}{-}} \FloatTok{0.5}
\NormalTok{migrate      }\OtherTok{\textless{}{-}} \FloatTok{0.01}

\NormalTok{output }\OtherTok{\textless{}{-}} \FunctionTok{simulate\_pop}\NormalTok{(}\DecValTok{600}\NormalTok{, }\DecValTok{0}\NormalTok{, }\FloatTok{0.01}\NormalTok{, }\FloatTok{0.2}\NormalTok{, }\FloatTok{0.001}\NormalTok{, max\_gen, K,}
\NormalTok{                       habitat\_prop, split\_ratio , split\_time,migrate)}

\CommentTok{\# overview plot}
\FunctionTok{plot\_split}\NormalTok{(output, }\AttributeTok{K =}\NormalTok{ K, }\AttributeTok{habitat\_prop =}\NormalTok{ habitat\_prop, }\AttributeTok{split\_time =}\NormalTok{ split\_time)}
\end{Highlighting}
\end{Shaded}

\pandocbounded{\includegraphics[keepaspectratio]{THEE_fragmentation_files/figure-latex/plotting-model-1.pdf}}

\begin{Shaded}
\begin{Highlighting}[]
\CommentTok{\# panels for habitat B and C}
\FunctionTok{plot\_habitats}\NormalTok{(output, }\AttributeTok{K =}\NormalTok{ K, }\AttributeTok{habitat\_prop =}\NormalTok{ habitat\_prop, }\AttributeTok{split\_time =}\NormalTok{ split\_time)}
\end{Highlighting}
\end{Shaded}

\pandocbounded{\includegraphics[keepaspectratio]{THEE_fragmentation_files/figure-latex/plotting-model-2.pdf}}

\begin{Shaded}
\begin{Highlighting}[]
\CommentTok{\# zoomed in on the first 150 generations}
\FunctionTok{plot\_split}\NormalTok{(output, }\AttributeTok{K =}\NormalTok{ K, }\AttributeTok{habitat\_prop =}\NormalTok{ habitat\_prop,}
           \AttributeTok{split\_time =}\NormalTok{ split\_time, }\AttributeTok{xlim\_max =} \DecValTok{150}\NormalTok{)}
\end{Highlighting}
\end{Shaded}

\pandocbounded{\includegraphics[keepaspectratio]{THEE_fragmentation_files/figure-latex/plotting-model-3.pdf}}

\begin{Shaded}
\begin{Highlighting}[]
\FunctionTok{plot\_habitats}\NormalTok{(output, }\AttributeTok{K =}\NormalTok{ K, }\AttributeTok{habitat\_prop =}\NormalTok{ habitat\_prop,}
              \AttributeTok{split\_time =}\NormalTok{ split\_time, }\AttributeTok{xlim\_max =} \DecValTok{150}\NormalTok{)}
\end{Highlighting}
\end{Shaded}

\pandocbounded{\includegraphics[keepaspectratio]{THEE_fragmentation_files/figure-latex/plotting-model-4.pdf}}

\section{Analyses and Results}\label{analyses-and-results}

A single simulation is one random realisation of the model: whether the
mutant arises, and whether it spreads before the wild type dies out,
depends on one run. To obtain interpretative results we repeat each
scenario many times and report the \textbf{rescue probability}, the
proportion of replicates in which rescue occured. The function
\texttt{run\_scenarios()} does this for any combination of
\texttt{split\_times}, \texttt{habitat\_props}, \texttt{split\_ratios}
and \texttt{migrates}, and \texttt{to\_long()} prepares the summary for
plotting.

All analyses use the same settings, and only the parameters named in
each analysis are varied. Note that migration is switched on in analyses
1 to 5 (migrate = 0.01) and only varied in the last analysis.

\subsubsection{\texorpdfstring{What \texttt{run\_scenarios()}
does}{What run\_scenarios() does}}\label{what-run_scenarios-does}

The function \texttt{run\_scenarios()} loops over every combination of
split time, habitat proportion, split ratio and migration rate, and each
combination is treated as one scenario. For each scenario it calls
\texttt{simulate\_pop()} repeatedly until 100 replicates have been made.
Runs in which the population went extinct (or reached twice its initial
size) before the split are discarded and drawn again, because habitats B
and C do not exist in these runs.

The helper function \texttt{to\_long()} reshapes the summary for
\texttt{ggplot2}: the columns \texttt{rescued\_B}, \texttt{rescued\_C}
and \texttt{rescued\_both} are combined into a single column, and the
standard error of a proportion, \(\sqrt{p(1-p)/n}\), is added and shown
as error bars in the plots.

\subsubsection{Important points for
interpretation}\label{important-points-for-interpretation}

Three features of this design affect how the results should be read.
First, runs that die out before the split are discarded, so all rescue
probabilities are conditional on the population still existing at the
split time. This matters most for late splits, where many runs without a
mutant have already gone extinct. Second, a habitat counts as rescued if
the number of mutants at the end of the run is at least
\texttt{rescue\_frac} times its carrying capacity. Because the
population fluctuates around K, this is a strict threshold. Third, with
100 replicates the standard error of a probability is up to about 0.05,
so small differences between scenarios should not be interpreted as real
effects.

\subsection{Overview of the analysis
code}\label{overview-of-the-analysis-code}

\begin{Shaded}
\begin{Highlighting}[]
\CommentTok{\# fixed parameters}
\NormalTok{init\_a           }\OtherTok{\textless{}{-}} \DecValTok{600}
\NormalTok{init\_A           }\OtherTok{\textless{}{-}} \DecValTok{0}
\NormalTok{m\_rate           }\OtherTok{\textless{}{-}} \FloatTok{0.001}
\NormalTok{max\_gen\_analysis }\OtherTok{\textless{}{-}} \DecValTok{1000}
\NormalTok{K\_analysis       }\OtherTok{\textless{}{-}} \DecValTok{1000}
\NormalTok{no\_replicates    }\OtherTok{\textless{}{-}} \DecValTok{100}    \CommentTok{\# more replicates give more precise probabilities}
\NormalTok{max\_attempts     }\OtherTok{\textless{}{-}} \DecValTok{5000}
\NormalTok{rescue\_frac      }\OtherTok{\textless{}{-}} \DecValTok{1}  
\NormalTok{migrate          }\OtherTok{\textless{}{-}} \FloatTok{0.01}    \CommentTok{\# rescue = A \textgreater{}= rescue\_frac * K of the habitat, changeable }

\CommentTok{\# core function}
\NormalTok{run\_scenarios }\OtherTok{\textless{}{-}} \ControlFlowTok{function}\NormalTok{(split\_times, habitat\_props, split\_ratios,}\AttributeTok{migrates =}\NormalTok{ migrate) \{}
\NormalTok{  results }\OtherTok{\textless{}{-}} \FunctionTok{list}\NormalTok{()}
  
  \ControlFlowTok{for}\NormalTok{ (st }\ControlFlowTok{in}\NormalTok{ split\_times) \{}
    \ControlFlowTok{for}\NormalTok{ (hp }\ControlFlowTok{in}\NormalTok{ habitat\_props) \{}
      \ControlFlowTok{for}\NormalTok{ (sr }\ControlFlowTok{in}\NormalTok{ split\_ratios) \{}
        \ControlFlowTok{for}\NormalTok{ (mg }\ControlFlowTok{in}\NormalTok{ migrates) \{  }
        
\NormalTok{        K\_B }\OtherTok{\textless{}{-}}\NormalTok{ K\_analysis }\SpecialCharTok{*}\NormalTok{ hp}
\NormalTok{        K\_C }\OtherTok{\textless{}{-}}\NormalTok{ K\_analysis }\SpecialCharTok{*}\NormalTok{ (}\DecValTok{1} \SpecialCharTok{{-}}\NormalTok{ hp)}
        
\NormalTok{        i }\OtherTok{\textless{}{-}} \DecValTok{0}
\NormalTok{        attempts }\OtherTok{\textless{}{-}} \DecValTok{0}
        
        \ControlFlowTok{repeat}\NormalTok{ \{}
\NormalTok{          attempts }\OtherTok{\textless{}{-}}\NormalTok{ attempts }\SpecialCharTok{+} \DecValTok{1}
          
          \ControlFlowTok{if}\NormalTok{ (attempts }\SpecialCharTok{\textgreater{}}\NormalTok{ max\_attempts) \{}
            \FunctionTok{warning}\NormalTok{(}\FunctionTok{paste}\NormalTok{(}
              \StringTok{"Too many failed runs for st ="}\NormalTok{, st,}
              \StringTok{"hp ="}\NormalTok{, hp,}
              \StringTok{"sr ="}\NormalTok{, sr,}
              \StringTok{"mig ="}\NormalTok{, mg}
\NormalTok{            ))}
            \ControlFlowTok{break}
\NormalTok{          \}}
          
\NormalTok{          one\_run }\OtherTok{\textless{}{-}} \FunctionTok{simulate\_pop}\NormalTok{(}
\NormalTok{            init\_a, init\_A,}
            \AttributeTok{decay\_rate   =} \FloatTok{0.1}\NormalTok{,}
            \AttributeTok{sel\_coeff    =} \FloatTok{0.2}\NormalTok{,}
            \AttributeTok{mut\_rate     =}\NormalTok{ m\_rate,}
            \AttributeTok{t\_max        =}\NormalTok{ max\_gen\_analysis,}
            \AttributeTok{K            =}\NormalTok{ K\_analysis,}
            \AttributeTok{habitat\_prop =}\NormalTok{ hp,}
            \AttributeTok{split\_ratio  =}\NormalTok{ sr,}
            \AttributeTok{split\_time   =}\NormalTok{ st,}
            \AttributeTok{migrate      =}\NormalTok{ mg}
\NormalTok{          )}
          
          \CommentTok{\# Skip runs that ended before the split time}
          \ControlFlowTok{if}\NormalTok{ (}\FunctionTok{nrow}\NormalTok{(one\_run}\SpecialCharTok{$}\NormalTok{B) }\SpecialCharTok{==} \DecValTok{0} \SpecialCharTok{||} \FunctionTok{nrow}\NormalTok{(one\_run}\SpecialCharTok{$}\NormalTok{C) }\SpecialCharTok{==} \DecValTok{0}\NormalTok{) }\ControlFlowTok{next}
          
\NormalTok{          i }\OtherTok{\textless{}{-}}\NormalTok{ i }\SpecialCharTok{+} \DecValTok{1}
          
\NormalTok{          reached\_B }\OtherTok{\textless{}{-}}\NormalTok{ one\_run}\SpecialCharTok{$}\NormalTok{B}\SpecialCharTok{$}\NormalTok{A }\SpecialCharTok{\textgreater{}=}\NormalTok{ rescue\_frac }\SpecialCharTok{*}\NormalTok{ K\_B}
\NormalTok{          reached\_C }\OtherTok{\textless{}{-}}\NormalTok{ one\_run}\SpecialCharTok{$}\NormalTok{C}\SpecialCharTok{$}\NormalTok{A }\SpecialCharTok{\textgreater{}=}\NormalTok{ rescue\_frac }\SpecialCharTok{*}\NormalTok{ K\_C}
          
\NormalTok{          rescued\_B }\OtherTok{\textless{}{-}}\NormalTok{ reached\_B[}\FunctionTok{length}\NormalTok{(reached\_B)]}
\NormalTok{          rescued\_C }\OtherTok{\textless{}{-}}\NormalTok{ reached\_C[}\FunctionTok{length}\NormalTok{(reached\_C)]}
          
\NormalTok{          results[[}\FunctionTok{length}\NormalTok{(results) }\SpecialCharTok{+} \DecValTok{1}\NormalTok{]] }\OtherTok{\textless{}{-}} \FunctionTok{data.frame}\NormalTok{(}
            \AttributeTok{st =}\NormalTok{ st,}
            \AttributeTok{hp =}\NormalTok{ hp,}
            \AttributeTok{sr =}\NormalTok{ sr,}
            \AttributeTok{mig =}\NormalTok{ mg,}
            \AttributeTok{replicate =}\NormalTok{ i,}
            \AttributeTok{rescued\_B =}\NormalTok{ rescued\_B,}
            \AttributeTok{rescued\_C =}\NormalTok{ rescued\_C,}
            \AttributeTok{rescued\_any =}\NormalTok{ rescued\_B }\SpecialCharTok{|}\NormalTok{ rescued\_C,}
            \AttributeTok{rescued\_both =}\NormalTok{ rescued\_B }\SpecialCharTok{\&}\NormalTok{ rescued\_C,}
            \AttributeTok{gen\_B =} \ControlFlowTok{if}\NormalTok{ (rescued\_B) }\FunctionTok{which}\NormalTok{(reached\_B)[}\DecValTok{1}\NormalTok{] }\ControlFlowTok{else} \ConstantTok{NA\_integer\_}\NormalTok{,}
            \AttributeTok{gen\_C =} \ControlFlowTok{if}\NormalTok{ (rescued\_C) }\FunctionTok{which}\NormalTok{(reached\_C)[}\DecValTok{1}\NormalTok{] }\ControlFlowTok{else} \ConstantTok{NA\_integer\_}
\NormalTok{          )}
          
          \ControlFlowTok{if}\NormalTok{ (i }\SpecialCharTok{\textgreater{}=}\NormalTok{ no\_replicates) }\ControlFlowTok{break}
\NormalTok{         \}}
\NormalTok{        \}}
\NormalTok{      \}}
\NormalTok{    \}}
\NormalTok{  \}}
  
  \CommentTok{\# Combine all simulation results}
  \ControlFlowTok{if}\NormalTok{ (}\FunctionTok{length}\NormalTok{(results) }\SpecialCharTok{==} \DecValTok{0}\NormalTok{) \{}
    \FunctionTok{stop}\NormalTok{(}\StringTok{"No valid simulation results were obtained."}\NormalTok{)}
\NormalTok{  \}}
  
\NormalTok{  data\_table }\OtherTok{\textless{}{-}} \FunctionTok{do.call}\NormalTok{(rbind, results)}

  \CommentTok{\# Rescue probability per scenario}
\NormalTok{  rescue\_prob }\OtherTok{\textless{}{-}}\NormalTok{ data\_table }\SpecialCharTok{\%\textgreater{}\%}
\NormalTok{    dplyr}\SpecialCharTok{::}\FunctionTok{group\_by}\NormalTok{(st, hp, sr, mig) }\SpecialCharTok{\%\textgreater{}\%}
\NormalTok{    dplyr}\SpecialCharTok{::}\FunctionTok{summarise}\NormalTok{(}
      \AttributeTok{n =}\NormalTok{ dplyr}\SpecialCharTok{::}\FunctionTok{n}\NormalTok{(),}
\NormalTok{      dplyr}\SpecialCharTok{::}\FunctionTok{across}\NormalTok{(}
        \FunctionTok{c}\NormalTok{(rescued\_B, rescued\_C, rescued\_any, rescued\_both),}
\NormalTok{        mean}
\NormalTok{      ),}
      \AttributeTok{.groups =} \StringTok{"drop"}
\NormalTok{    )}
  
  \FunctionTok{return}\NormalTok{(}\FunctionTok{list}\NormalTok{(}
    \AttributeTok{data =}\NormalTok{ data\_table,}
    \AttributeTok{prob =}\NormalTok{ rescue\_prob}
\NormalTok{  ))}
\NormalTok{\}}


\CommentTok{\# {-}{-}{-}{-} helperfunction : long format with standard error for plotting {-}{-}{-}{-}}

\NormalTok{to\_long }\OtherTok{\textless{}{-}} \ControlFlowTok{function}\NormalTok{(rescue\_prob) \{}
\NormalTok{  rescue\_prob }\SpecialCharTok{\%\textgreater{}\%}         \CommentTok{\#\%\textgreater{}\% is the pipe operator. It takes whatever is on its left and passes it as the first argument of the function on its right.}
\NormalTok{    tidyr}\SpecialCharTok{::}\FunctionTok{pivot\_longer}\NormalTok{(   }\CommentTok{\#doing this for ggplot 2 }
      \AttributeTok{cols =} \FunctionTok{c}\NormalTok{(rescued\_B, rescued\_C, rescued\_both),}
      \AttributeTok{names\_to =} \StringTok{"population"}\NormalTok{,}
      \AttributeTok{values\_to =} \StringTok{"rescue\_probability"}   
\NormalTok{    ) }\SpecialCharTok{\%\textgreater{}\%}
\NormalTok{    dplyr}\SpecialCharTok{::}\FunctionTok{mutate}\NormalTok{(}
      \CommentTok{\#standard error of a proportion}
      \AttributeTok{se =} \FunctionTok{sqrt}\NormalTok{(rescue\_probability }\SpecialCharTok{*}\NormalTok{ (}\DecValTok{1} \SpecialCharTok{{-}}\NormalTok{ rescue\_probability) }\SpecialCharTok{/}\NormalTok{ n),}
      \AttributeTok{population =}\NormalTok{ dplyr}\SpecialCharTok{::}\FunctionTok{recode}\NormalTok{(}
\NormalTok{        population,}
        \AttributeTok{rescued\_B    =} \StringTok{"Population B"}\NormalTok{,}
        \AttributeTok{rescued\_C    =} \StringTok{"Population C"}\NormalTok{,}
        \AttributeTok{rescued\_both =} \StringTok{"Both populations"}
\NormalTok{      )}
\NormalTok{    )}
\NormalTok{\}}
\end{Highlighting}
\end{Shaded}

\begin{Shaded}
\begin{Highlighting}[]
\CommentTok{\# =======================================================================}
\CommentTok{\# ANALYSIS 1: split time varies, habitat proportion fixed}
\CommentTok{\# =======================================================================}
\NormalTok{res1 }\OtherTok{\textless{}{-}} \FunctionTok{run\_scenarios}\NormalTok{(}
  \AttributeTok{split\_times   =} \FunctionTok{c}\NormalTok{(}\DecValTok{10}\NormalTok{, }\DecValTok{20}\NormalTok{, }\DecValTok{30}\NormalTok{, }\DecValTok{50}\NormalTok{, }\DecValTok{75}\NormalTok{, }\DecValTok{100}\NormalTok{),}
  \AttributeTok{habitat\_props =} \FloatTok{0.5}\NormalTok{,}
  \AttributeTok{split\_ratios =} \FloatTok{0.5}
\NormalTok{)}
\NormalTok{long1 }\OtherTok{\textless{}{-}} \FunctionTok{to\_long}\NormalTok{(res1}\SpecialCharTok{$}\NormalTok{prob)}
\FunctionTok{print}\NormalTok{(res1}\SpecialCharTok{$}\NormalTok{prob)}

\FunctionTok{ggplot}\NormalTok{(long1, }\FunctionTok{aes}\NormalTok{(}\AttributeTok{x =}\NormalTok{ st, }\AttributeTok{y =}\NormalTok{ rescue\_probability,}
                  \AttributeTok{color =}\NormalTok{ population, }\AttributeTok{group =}\NormalTok{ population)) }\SpecialCharTok{+}
  \FunctionTok{geom\_line}\NormalTok{(}\AttributeTok{linewidth =} \DecValTok{1}\NormalTok{) }\SpecialCharTok{+}
  \FunctionTok{geom\_point}\NormalTok{(}\AttributeTok{size =} \DecValTok{3}\NormalTok{) }\SpecialCharTok{+}
  \FunctionTok{geom\_errorbar}\NormalTok{(}\FunctionTok{aes}\NormalTok{(}\AttributeTok{ymin =} \FunctionTok{pmax}\NormalTok{(}\DecValTok{0}\NormalTok{, rescue\_probability }\SpecialCharTok{{-}}\NormalTok{ se),}
                    \AttributeTok{ymax =} \FunctionTok{pmin}\NormalTok{(}\DecValTok{1}\NormalTok{, rescue\_probability }\SpecialCharTok{+}\NormalTok{ se)),}
                \AttributeTok{width =} \DecValTok{2}\NormalTok{) }\SpecialCharTok{+}
  \FunctionTok{scale\_y\_continuous}\NormalTok{(}\AttributeTok{limits =} \FunctionTok{c}\NormalTok{(}\DecValTok{0}\NormalTok{, }\DecValTok{1}\NormalTok{), }\AttributeTok{breaks =} \FunctionTok{seq}\NormalTok{(}\DecValTok{0}\NormalTok{, }\DecValTok{1}\NormalTok{, }\FloatTok{0.1}\NormalTok{)) }\SpecialCharTok{+}
  \FunctionTok{labs}\NormalTok{(}\AttributeTok{title =} \StringTok{"Rescue prob. for split time habitat proportion = 0.5"}\NormalTok{,}
       \AttributeTok{x =} \StringTok{"Split time"}\NormalTok{, }\AttributeTok{y =} \StringTok{"Rescue probability"}\NormalTok{, }\AttributeTok{color =} \StringTok{"Population"}\NormalTok{) }\SpecialCharTok{+}
  \FunctionTok{theme\_minimal}\NormalTok{(}\AttributeTok{base\_size =} \DecValTok{16}\NormalTok{)}
\end{Highlighting}
\end{Shaded}

\pandocbounded{\includegraphics[keepaspectratio]{THEE_fragmentation_files/figure-latex/analyses-1-1.pdf}}
\#\# Timing of fragmentation This plot shows the proportion of 100
replicate simulations in which the mutant reached the rescue threshold
(100\% of K), depending on the generation at which the habitat was split
into two equal halves. Rescue probability tended to increase with later
splits, from about 0.24 (B) and 0.23 (C) at generation 10 to about 0.43
and 0.41 at generation 100, with a dip at generation 30. Habitats B and
C behaved very similarly, as expected for identical halves. The
probability that both were rescued was lower (about 0.06 to 0.17) and
stayed roughly flat from generation 50 onwards.

\begin{Shaded}
\begin{Highlighting}[]
\CommentTok{\# =======================================================================}
\CommentTok{\# ANALYSIS 2: habitat proportion varies, split time fixed}
\CommentTok{\# =======================================================================}
\NormalTok{res2 }\OtherTok{\textless{}{-}} \FunctionTok{run\_scenarios}\NormalTok{(}
  \AttributeTok{split\_times   =} \DecValTok{30}\NormalTok{,}
  \AttributeTok{habitat\_props =} \FunctionTok{c}\NormalTok{(}\FloatTok{0.1}\NormalTok{, }\FloatTok{0.2}\NormalTok{, }\FloatTok{0.3}\NormalTok{, }\FloatTok{0.4}\NormalTok{, }\FloatTok{0.5}\NormalTok{),}
  \AttributeTok{split\_ratios =} \FloatTok{0.5}
\NormalTok{)}
\NormalTok{long2 }\OtherTok{\textless{}{-}} \FunctionTok{to\_long}\NormalTok{(res2}\SpecialCharTok{$}\NormalTok{prob)}
\FunctionTok{print}\NormalTok{(res2}\SpecialCharTok{$}\NormalTok{prob)}

\FunctionTok{ggplot}\NormalTok{(long2, }\FunctionTok{aes}\NormalTok{(}\AttributeTok{x =}\NormalTok{ hp, }\AttributeTok{y =}\NormalTok{ rescue\_probability,}
                  \AttributeTok{color =}\NormalTok{ population, }\AttributeTok{group =}\NormalTok{ population)) }\SpecialCharTok{+}
  \FunctionTok{geom\_line}\NormalTok{(}\AttributeTok{linewidth =} \DecValTok{1}\NormalTok{) }\SpecialCharTok{+}
  \FunctionTok{geom\_point}\NormalTok{(}\AttributeTok{size =} \DecValTok{3}\NormalTok{) }\SpecialCharTok{+}
  \FunctionTok{geom\_errorbar}\NormalTok{(}\FunctionTok{aes}\NormalTok{(}\AttributeTok{ymin =} \FunctionTok{pmax}\NormalTok{(}\DecValTok{0}\NormalTok{, rescue\_probability }\SpecialCharTok{{-}}\NormalTok{ se),}
                    \AttributeTok{ymax =} \FunctionTok{pmin}\NormalTok{(}\DecValTok{1}\NormalTok{, rescue\_probability }\SpecialCharTok{+}\NormalTok{ se)),}
                \AttributeTok{width =} \FloatTok{0.02}\NormalTok{) }\SpecialCharTok{+}
  \FunctionTok{scale\_x\_continuous}\NormalTok{(}\AttributeTok{breaks =} \FunctionTok{seq}\NormalTok{(}\FloatTok{0.1}\NormalTok{, }\FloatTok{0.9}\NormalTok{, }\FloatTok{0.1}\NormalTok{)) }\SpecialCharTok{+}
  \FunctionTok{scale\_y\_continuous}\NormalTok{(}\AttributeTok{limits =} \FunctionTok{c}\NormalTok{(}\DecValTok{0}\NormalTok{, }\DecValTok{1}\NormalTok{), }\AttributeTok{breaks =} \FunctionTok{seq}\NormalTok{(}\DecValTok{0}\NormalTok{, }\DecValTok{1}\NormalTok{, }\FloatTok{0.1}\NormalTok{)) }\SpecialCharTok{+}
  \FunctionTok{labs}\NormalTok{(}\AttributeTok{title =} \StringTok{"Rescue Prob. for habitat proportions (split time = 30)"}\NormalTok{,}
       \AttributeTok{x =} \StringTok{"Habitat proportion (B)"}\NormalTok{, }\AttributeTok{y =} \StringTok{"Rescue probability"}\NormalTok{,}
       \AttributeTok{color =} \StringTok{"Population"}\NormalTok{) }\SpecialCharTok{+}
  \FunctionTok{theme\_minimal}\NormalTok{(}\AttributeTok{base\_size =} \DecValTok{16}\NormalTok{)}
\end{Highlighting}
\end{Shaded}

\pandocbounded{\includegraphics[keepaspectratio]{THEE_fragmentation_files/figure-latex/analyses-2-1.pdf}}
\#\# Habitat proportion This Figure shows how rescue probability changes
when the habitat is split differently at generation 30. The x-axis gives
the share of the original habitat and also of K that goes to habitat B,
with habitat C receiving the remainder, while individuals were still
divided evenly. The smaller habitat was rescued more often. At a habitat
proportion of 0.1, habitat B (10\% of K) was rescued in 39\% of the 100
replicates and habitat C (90\% of K) in only 5\%. As B grew towards
equal size, its rescue probability fell (0.18 at 0.5) and that of C rose
(0.26 at 0.5). The probability that both habitats were rescued stayed
low at 0.05 to 0.10.

\begin{Shaded}
\begin{Highlighting}[]
\CommentTok{\# =======================================================================}
\CommentTok{\# ANALYSIS 3: split ratio varies, habitat proportion fixed}
\CommentTok{\# =======================================================================}

\NormalTok{res3 }\OtherTok{\textless{}{-}} \FunctionTok{run\_scenarios}\NormalTok{(}
  \AttributeTok{split\_times   =} \DecValTok{30}\NormalTok{,  }\CommentTok{\# fix split time}
  \AttributeTok{habitat\_props =} \FloatTok{0.5}\NormalTok{,}
  \AttributeTok{split\_ratios  =} \FunctionTok{c}\NormalTok{(}\FloatTok{0.1}\NormalTok{, }\FloatTok{0.2}\NormalTok{, }\FloatTok{0.3}\NormalTok{, }\FloatTok{0.4}\NormalTok{, }\FloatTok{0.5}\NormalTok{)}
\NormalTok{)}

\NormalTok{long3 }\OtherTok{\textless{}{-}} \FunctionTok{to\_long}\NormalTok{(res3}\SpecialCharTok{$}\NormalTok{prob)}

\FunctionTok{print}\NormalTok{(res3}\SpecialCharTok{$}\NormalTok{prob)}
\end{Highlighting}
\end{Shaded}

\begin{verbatim}
## # A tibble: 5 x 9
##      st    hp    sr   mig     n rescued_B rescued_C rescued_any rescued_both
##   <dbl> <dbl> <dbl> <dbl> <int>     <dbl>     <dbl>       <dbl>        <dbl>
## 1    30   0.5   0.1  0.01   100      0.17      0.24        0.34         0.07
## 2    30   0.5   0.2  0.01   100      0.19      0.22        0.33         0.08
## 3    30   0.5   0.3  0.01   100      0.32      0.33        0.49         0.16
## 4    30   0.5   0.4  0.01   100      0.24      0.24        0.34         0.14
## 5    30   0.5   0.5  0.01   100      0.22      0.19        0.3          0.11
\end{verbatim}

\begin{Shaded}
\begin{Highlighting}[]
\FunctionTok{ggplot}\NormalTok{(long3, }\FunctionTok{aes}\NormalTok{(}\AttributeTok{x =}\NormalTok{ sr, }\AttributeTok{y =}\NormalTok{ rescue\_probability,}
                  \AttributeTok{color =}\NormalTok{ population, }\AttributeTok{group =}\NormalTok{ population)) }\SpecialCharTok{+}
  \FunctionTok{geom\_line}\NormalTok{(}\AttributeTok{linewidth =} \DecValTok{1}\NormalTok{) }\SpecialCharTok{+}
  \FunctionTok{geom\_point}\NormalTok{(}\AttributeTok{size =} \DecValTok{3}\NormalTok{) }\SpecialCharTok{+}
  \FunctionTok{geom\_errorbar}\NormalTok{(}
    \FunctionTok{aes}\NormalTok{(}
      \AttributeTok{ymin =} \FunctionTok{pmax}\NormalTok{(}\DecValTok{0}\NormalTok{, rescue\_probability }\SpecialCharTok{{-}}\NormalTok{ se),}
      \AttributeTok{ymax =} \FunctionTok{pmin}\NormalTok{(}\DecValTok{1}\NormalTok{, rescue\_probability }\SpecialCharTok{+}\NormalTok{ se)}
\NormalTok{    ),}
    \AttributeTok{width =} \FloatTok{0.02}
\NormalTok{  ) }\SpecialCharTok{+}
  \FunctionTok{scale\_x\_continuous}\NormalTok{(}\AttributeTok{breaks =} \FunctionTok{seq}\NormalTok{(}\FloatTok{0.2}\NormalTok{, }\FloatTok{0.8}\NormalTok{, }\FloatTok{0.1}\NormalTok{)) }\SpecialCharTok{+}
  \FunctionTok{scale\_y\_continuous}\NormalTok{(}
    \AttributeTok{limits =} \FunctionTok{c}\NormalTok{(}\DecValTok{0}\NormalTok{, }\DecValTok{1}\NormalTok{),}
    \AttributeTok{breaks =} \FunctionTok{seq}\NormalTok{(}\DecValTok{0}\NormalTok{, }\DecValTok{1}\NormalTok{, }\FloatTok{0.1}\NormalTok{)}
\NormalTok{  ) }\SpecialCharTok{+}
  \FunctionTok{labs}\NormalTok{(}
    \AttributeTok{title =} \StringTok{"Rescue probability vs. split ratio (habitat proportion = 0.5)"}\NormalTok{,}
    \AttributeTok{x =} \StringTok{"Split ratio (sr)"}\NormalTok{,}
    \AttributeTok{y =} \StringTok{"Rescue probability"}\NormalTok{,}
    \AttributeTok{color =} \StringTok{"Population"}
\NormalTok{  ) }\SpecialCharTok{+}
  \FunctionTok{theme\_minimal}\NormalTok{(}\AttributeTok{base\_size =} \DecValTok{16}\NormalTok{)}
\end{Highlighting}
\end{Shaded}

\pandocbounded{\includegraphics[keepaspectratio]{THEE_fragmentation_files/figure-latex/analyses-3-1.pdf}}
\#\# Distribution of individuals the third analyses shows the effect of
how individuals are distributed at the split, with habitats of equal
size split at generation 30. The x-axis gives the probability (split
ratio) that an individual is assigned to habitat B, with the remaining
individuals going to C. Rescue probability showed no consistent trend
across split ratios. Habitat B was rescued in 17\% to 32\% of the 100
replicates and habitat C in 19\% to 33\%, and the two habitats did not
differ clearly at any ratio. Both curves peaked at a split ratio of 0.3.
The probability that both habitats were rescued was lower at more uneven
splits.

\begin{Shaded}
\begin{Highlighting}[]
\CommentTok{\# =======================================================================}
\CommentTok{\# ANALYSIS 4: both vary (full grid)}
\CommentTok{\# =======================================================================}
\NormalTok{res4 }\OtherTok{\textless{}{-}} \FunctionTok{run\_scenarios}\NormalTok{(}
  \AttributeTok{split\_times   =} \FunctionTok{c}\NormalTok{(}\DecValTok{10}\NormalTok{, }\DecValTok{30}\NormalTok{, }\DecValTok{50}\NormalTok{, }\DecValTok{100}\NormalTok{),}
  \AttributeTok{habitat\_props =} \FunctionTok{c}\NormalTok{(}\FloatTok{0.1}\NormalTok{, }\FloatTok{0.25}\NormalTok{, }\FloatTok{0.5}\NormalTok{),}
  \AttributeTok{split\_ratios =} \FloatTok{0.5}
\NormalTok{)}
\NormalTok{long4 }\OtherTok{\textless{}{-}} \FunctionTok{to\_long}\NormalTok{(res4}\SpecialCharTok{$}\NormalTok{prob)}
\FunctionTok{print}\NormalTok{(res4}\SpecialCharTok{$}\NormalTok{prob)}
\end{Highlighting}
\end{Shaded}

\begin{verbatim}
## # A tibble: 12 x 9
##       st    hp    sr   mig     n rescued_B rescued_C rescued_any rescued_both
##    <dbl> <dbl> <dbl> <dbl> <int>     <dbl>     <dbl>       <dbl>        <dbl>
##  1    10  0.1    0.5  0.01   100      0.44      0.05        0.44         0.05
##  2    10  0.25   0.5  0.01   100      0.32      0.06        0.32         0.06
##  3    10  0.5    0.5  0.01   100      0.17      0.13        0.22         0.08
##  4    30  0.1    0.5  0.01   100      0.5       0.09        0.5          0.09
##  5    30  0.25   0.5  0.01   100      0.39      0.1         0.4          0.09
##  6    30  0.5    0.5  0.01   100      0.23      0.17        0.3          0.1 
##  7    50  0.1    0.5  0.01   100      0.74      0.06        0.74         0.06
##  8    50  0.25   0.5  0.01   100      0.71      0.17        0.73         0.15
##  9    50  0.5    0.5  0.01   100      0.46      0.43        0.64         0.25
## 10   100  0.1    0.5  0.01   100      1         0.14        1            0.14
## 11   100  0.25   0.5  0.01   100      0.85      0.17        0.87         0.15
## 12   100  0.5    0.5  0.01   100      0.53      0.46        0.78         0.21
\end{verbatim}

\begin{Shaded}
\begin{Highlighting}[]
\CommentTok{\# heatmap: one panel per population}
\CommentTok{\# label showing the proportion of BOTH habitats on two lines, e.g. "B 0.1" over "C 0.9"}
\NormalTok{hp\_values }\OtherTok{\textless{}{-}} \FunctionTok{sort}\NormalTok{(}\FunctionTok{unique}\NormalTok{(long4}\SpecialCharTok{$}\NormalTok{hp))}
\NormalTok{long4 }\OtherTok{\textless{}{-}}\NormalTok{ long4 }\SpecialCharTok{\%\textgreater{}\%}
  \FunctionTok{mutate}\NormalTok{(}\AttributeTok{hp\_label =} \FunctionTok{factor}\NormalTok{(hp, }\AttributeTok{levels =}\NormalTok{ hp\_values,}
                           \AttributeTok{labels =} \FunctionTok{paste0}\NormalTok{(}\StringTok{"B "}\NormalTok{, hp\_values, }\StringTok{"}\SpecialCharTok{\textbackslash{}n}\StringTok{C "}\NormalTok{, }\DecValTok{1} \SpecialCharTok{{-}}\NormalTok{ hp\_values)))}

\FunctionTok{ggplot}\NormalTok{(long4, }\FunctionTok{aes}\NormalTok{(}\AttributeTok{x =}\NormalTok{ hp\_label, }\AttributeTok{y =} \FunctionTok{factor}\NormalTok{(st), }\AttributeTok{fill =}\NormalTok{ rescue\_probability)) }\SpecialCharTok{+}
  \FunctionTok{geom\_tile}\NormalTok{(}\AttributeTok{color =} \StringTok{"white"}\NormalTok{) }\SpecialCharTok{+}
  \FunctionTok{geom\_text}\NormalTok{(}\FunctionTok{aes}\NormalTok{(}\AttributeTok{label =} \FunctionTok{sprintf}\NormalTok{(}\StringTok{"\%.2f"}\NormalTok{, rescue\_probability)), }\AttributeTok{size =} \FloatTok{3.5}\NormalTok{) }\SpecialCharTok{+}
  \FunctionTok{facet\_wrap}\NormalTok{(}\SpecialCharTok{\textasciitilde{}}\NormalTok{ population) }\SpecialCharTok{+}
  \FunctionTok{scale\_fill\_viridis\_c}\NormalTok{(}\AttributeTok{limits =} \FunctionTok{c}\NormalTok{(}\DecValTok{0}\NormalTok{, }\DecValTok{1}\NormalTok{), }\AttributeTok{option =} \StringTok{"C"}\NormalTok{) }\SpecialCharTok{+}
  \FunctionTok{labs}\NormalTok{(}\AttributeTok{title =} \StringTok{"Rescue probability by split time and habitat proportion"}\NormalTok{,}
       \AttributeTok{x =} \StringTok{"Habitat proportion"}\NormalTok{, }\AttributeTok{y =} \StringTok{"Split time"}\NormalTok{,}
       \AttributeTok{fill =} \StringTok{"Rescue}\SpecialCharTok{\textbackslash{}n}\StringTok{probability"}\NormalTok{) }\SpecialCharTok{+}
  \FunctionTok{theme\_minimal}\NormalTok{(}\AttributeTok{base\_size=} \DecValTok{16}\NormalTok{)}
\end{Highlighting}
\end{Shaded}

\pandocbounded{\includegraphics[keepaspectratio]{THEE_fragmentation_files/figure-latex/analyses-4-1.pdf}}
\#\# Split time x habitat proportion This heatmap combines both factors
for habitat proportions of 0.1, 0.25 and 0.5. Rescue of habitat B
increased with split time at every proportion. It was highest when B was
smallest: at a proportion of 0.1 it rose from 0.44 to 1.00, compared
with 0.17 to 0.53 at 0.5. Habitat C, the larger habitat in unequal
splits, was rarely rescued (at most 0.17). Only at equal size did its
rescue probability increase with split time (0.13 to 0.46). Rescue of
both habitats stayed low (0.05 to 0.25). The smaller habitat was thus
rescued more often than the larger one. This difference grew with later
splits.

\begin{Shaded}
\begin{Highlighting}[]
\CommentTok{\# =======================================================================}
\CommentTok{\# ANALYSIS 5: split time and split ratio vary (full grid)}
\CommentTok{\# =======================================================================}

\NormalTok{res5 }\OtherTok{\textless{}{-}} \FunctionTok{run\_scenarios}\NormalTok{(}
  \AttributeTok{split\_times   =} \FunctionTok{c}\NormalTok{(}\DecValTok{10}\NormalTok{, }\DecValTok{30}\NormalTok{, }\DecValTok{50}\NormalTok{, }\DecValTok{100}\NormalTok{),}
  \AttributeTok{habitat\_props =} \FloatTok{0.5}\NormalTok{,}
  \AttributeTok{split\_ratios  =} \FunctionTok{c}\NormalTok{(}\FloatTok{0.1}\NormalTok{, }\FloatTok{0.25}\NormalTok{, }\FloatTok{0.5}\NormalTok{)}
\NormalTok{)}

\NormalTok{long5 }\OtherTok{\textless{}{-}} \FunctionTok{to\_long}\NormalTok{(res5}\SpecialCharTok{$}\NormalTok{prob)}

\FunctionTok{print}\NormalTok{(res5}\SpecialCharTok{$}\NormalTok{prob)}
\end{Highlighting}
\end{Shaded}

\begin{verbatim}
## # A tibble: 12 x 9
##       st    hp    sr   mig     n rescued_B rescued_C rescued_any rescued_both
##    <dbl> <dbl> <dbl> <dbl> <int>     <dbl>     <dbl>       <dbl>        <dbl>
##  1    10   0.5  0.1   0.01   100      0.23      0.24        0.38         0.09
##  2    10   0.5  0.25  0.01   100      0.22      0.24        0.37         0.09
##  3    10   0.5  0.5   0.01   100      0.16      0.18        0.29         0.05
##  4    30   0.5  0.1   0.01   100      0.2       0.27        0.36         0.11
##  5    30   0.5  0.25  0.01   100      0.27      0.16        0.32         0.11
##  6    30   0.5  0.5   0.01   100      0.23      0.19        0.34         0.08
##  7    50   0.5  0.1   0.01   100      0.34      0.37        0.52         0.19
##  8    50   0.5  0.25  0.01   100      0.35      0.36        0.52         0.19
##  9    50   0.5  0.5   0.01   100      0.31      0.37        0.51         0.17
## 10   100   0.5  0.1   0.01   100      0.39      0.51        0.7          0.2 
## 11   100   0.5  0.25  0.01   100      0.43      0.45        0.67         0.21
## 12   100   0.5  0.5   0.01   100      0.48      0.46        0.71         0.23
\end{verbatim}

\begin{Shaded}
\begin{Highlighting}[]
\CommentTok{\# Heatmap: one panel per population}
\FunctionTok{ggplot}\NormalTok{(long5, }\FunctionTok{aes}\NormalTok{(}\AttributeTok{x =} \FunctionTok{factor}\NormalTok{(sr), }\AttributeTok{y =} \FunctionTok{factor}\NormalTok{(st),}
                  \AttributeTok{fill =}\NormalTok{ rescue\_probability)) }\SpecialCharTok{+}
  \FunctionTok{geom\_tile}\NormalTok{(}\AttributeTok{color =} \StringTok{"white"}\NormalTok{) }\SpecialCharTok{+}
  \FunctionTok{geom\_text}\NormalTok{(}\FunctionTok{aes}\NormalTok{(}\AttributeTok{label =} \FunctionTok{sprintf}\NormalTok{(}\StringTok{"\%.2f"}\NormalTok{, rescue\_probability)),}
            \AttributeTok{size =} \FloatTok{3.5}\NormalTok{) }\SpecialCharTok{+}
  \FunctionTok{facet\_wrap}\NormalTok{(}\SpecialCharTok{\textasciitilde{}}\NormalTok{ population) }\SpecialCharTok{+}
  \FunctionTok{scale\_fill\_viridis\_c}\NormalTok{(}\AttributeTok{limits =} \FunctionTok{c}\NormalTok{(}\DecValTok{0}\NormalTok{, }\DecValTok{1}\NormalTok{), }\AttributeTok{option =} \StringTok{"C"}\NormalTok{) }\SpecialCharTok{+}
  \FunctionTok{labs}\NormalTok{(}
    \AttributeTok{title =} \StringTok{"Rescue probability by split time and split ratio"}\NormalTok{,}
    \AttributeTok{x =} \StringTok{"Split ratio (sr)"}\NormalTok{,}
    \AttributeTok{y =} \StringTok{"Split time (st)"}\NormalTok{,}
    \AttributeTok{fill =} \StringTok{"Rescue}\SpecialCharTok{\textbackslash{}n}\StringTok{probability"}
\NormalTok{  ) }\SpecialCharTok{+}
  \FunctionTok{theme\_minimal}\NormalTok{(}\AttributeTok{base\_size =} \DecValTok{16}\NormalTok{)}
\end{Highlighting}
\end{Shaded}

\pandocbounded{\includegraphics[keepaspectratio]{THEE_fragmentation_files/figure-latex/analyses-5-1.pdf}}
\#\# Split time x split ratio This Plot combines split time with split
ratios of 0.1, 0.25 and 0.5. Rescue probability increased with split
time in both habitats. For habitat B it rose from about 0.2 at
generation 10 to about 0.4 to 0.5 at generation 100. Habitat C showed
the same pattern. Rescue of both habitats also increased, from 0.05 to
0.09 at generation 10 to about 0.2 at generation 100. The split ratio
had little effect. Differences between ratios were small and showed no
clear direction.

\begin{Shaded}
\begin{Highlighting}[]
\CommentTok{\# =======================================================================}
\CommentTok{\# ANALYSIS 6: migration rate and habitat proportion vary (full grid)}
\CommentTok{\# =======================================================================}
\NormalTok{res7 }\OtherTok{\textless{}{-}} \FunctionTok{run\_scenarios}\NormalTok{(}
  \AttributeTok{split\_times   =} \DecValTok{30}\NormalTok{,}
  \AttributeTok{habitat\_props =} \FunctionTok{c}\NormalTok{(}\FloatTok{0.1}\NormalTok{, }\FloatTok{0.25}\NormalTok{, }\FloatTok{0.5}\NormalTok{),}
  \AttributeTok{split\_ratios  =} \FloatTok{0.5}\NormalTok{,}
  \AttributeTok{migrates      =} \FunctionTok{c}\NormalTok{(}\DecValTok{0}\NormalTok{, }\FloatTok{0.01}\NormalTok{, }\FloatTok{0.05}\NormalTok{, }\FloatTok{0.1}\NormalTok{)}
\NormalTok{)}
\NormalTok{long7 }\OtherTok{\textless{}{-}} \FunctionTok{to\_long}\NormalTok{(res7}\SpecialCharTok{$}\NormalTok{prob)}
\FunctionTok{print}\NormalTok{(res7}\SpecialCharTok{$}\NormalTok{prob)}
\end{Highlighting}
\end{Shaded}

\begin{verbatim}
## # A tibble: 12 x 9
##       st    hp    sr   mig     n rescued_B rescued_C rescued_any rescued_both
##    <dbl> <dbl> <dbl> <dbl> <int>     <dbl>     <dbl>       <dbl>        <dbl>
##  1    30  0.1    0.5  0      100      0.23      0.24        0.38         0.09
##  2    30  0.1    0.5  0.01   100      0.51      0.06        0.51         0.06
##  3    30  0.1    0.5  0.05   100      0.44      0           0.44         0   
##  4    30  0.1    0.5  0.1    100      0.48      0           0.48         0   
##  5    30  0.25   0.5  0      100      0.21      0.16        0.32         0.05
##  6    30  0.25   0.5  0.01   100      0.38      0.12        0.39         0.11
##  7    30  0.25   0.5  0.05   100      0.47      0           0.47         0   
##  8    30  0.25   0.5  0.1    100      0.47      0           0.47         0   
##  9    30  0.5    0.5  0      100      0.21      0.21        0.31         0.11
## 10    30  0.5    0.5  0.01   100      0.22      0.21        0.35         0.08
## 11    30  0.5    0.5  0.05   100      0.3       0.21        0.33         0.18
## 12    30  0.5    0.5  0.1    100      0.19      0.23        0.28         0.14
\end{verbatim}

\begin{Shaded}
\begin{Highlighting}[]
\CommentTok{\# label showing the proportion of BOTH habitats on two lines}
\NormalTok{hp\_values7 }\OtherTok{\textless{}{-}} \FunctionTok{sort}\NormalTok{(}\FunctionTok{unique}\NormalTok{(long7}\SpecialCharTok{$}\NormalTok{hp))}
\NormalTok{long7 }\OtherTok{\textless{}{-}}\NormalTok{ long7 }\SpecialCharTok{\%\textgreater{}\%}
  \FunctionTok{mutate}\NormalTok{(}\AttributeTok{hp\_label =} \FunctionTok{factor}\NormalTok{(hp, }\AttributeTok{levels =}\NormalTok{ hp\_values7,}
                           \AttributeTok{labels =} \FunctionTok{paste0}\NormalTok{(}\StringTok{"B "}\NormalTok{, hp\_values7, }\StringTok{"}\SpecialCharTok{\textbackslash{}n}\StringTok{C "}\NormalTok{, }\DecValTok{1} \SpecialCharTok{{-}}\NormalTok{ hp\_values7)))}

\FunctionTok{ggplot}\NormalTok{(long7, }\FunctionTok{aes}\NormalTok{(}\AttributeTok{x =}\NormalTok{ hp\_label, }\AttributeTok{y =} \FunctionTok{factor}\NormalTok{(mig), }\AttributeTok{fill =}\NormalTok{ rescue\_probability)) }\SpecialCharTok{+}
  \FunctionTok{geom\_tile}\NormalTok{(}\AttributeTok{color =} \StringTok{"white"}\NormalTok{) }\SpecialCharTok{+}
  \FunctionTok{geom\_text}\NormalTok{(}\FunctionTok{aes}\NormalTok{(}\AttributeTok{label =} \FunctionTok{sprintf}\NormalTok{(}\StringTok{"\%.2f"}\NormalTok{, rescue\_probability)), }\AttributeTok{size =} \FloatTok{3.5}\NormalTok{) }\SpecialCharTok{+}
  \FunctionTok{facet\_wrap}\NormalTok{(}\SpecialCharTok{\textasciitilde{}}\NormalTok{ population) }\SpecialCharTok{+}
  \FunctionTok{scale\_fill\_viridis\_c}\NormalTok{(}\AttributeTok{limits =} \FunctionTok{c}\NormalTok{(}\DecValTok{0}\NormalTok{, }\DecValTok{1}\NormalTok{), }\AttributeTok{option =} \StringTok{"C"}\NormalTok{) }\SpecialCharTok{+}
  \FunctionTok{labs}\NormalTok{(}\AttributeTok{title =} \StringTok{"Rescue proba. by migration rate and habitat propo."}\NormalTok{,}
       \AttributeTok{x =} \StringTok{"Habitat proportion"}\NormalTok{, }\AttributeTok{y =} \StringTok{"Migration rate"}\NormalTok{,}
       \AttributeTok{fill =} \StringTok{"Rescue}\SpecialCharTok{\textbackslash{}n}\StringTok{probability"}\NormalTok{) }\SpecialCharTok{+}
  \FunctionTok{theme\_minimal}\NormalTok{(}\AttributeTok{base\_size =} \DecValTok{16}\NormalTok{)}
\end{Highlighting}
\end{Shaded}

\pandocbounded{\includegraphics[keepaspectratio]{THEE_fragmentation_files/figure-latex/analysis-6-migration-habitat-1.pdf}}
\#\#\# Migration x habitat proportion The last analyses combines the
migration rate with habitat proportions of 0.1, 0.25 and 0.5. Without
migration, the habitats were rescued about equally often, at every
proportion. Even the smallest habitat had no advantage. With migration,
unequal habitats behaved very differently. The smaller habitat B was
rescued more often. The larger habitat C was rarely rescued at 0.01, at
most 0.11 and never at 0.05 or higher. Rescue of both habitats dropped
to zero accordingly.

\section{Discussion}\label{discussion}

\section*{Appendix}\label{appendix}
\addcontentsline{toc}{section}{Appendix}

\textbf{Less common functions used}

\begin{itemize}
\tightlist
\item
  \texttt{set.seed(57)}: initialises R's random number generator. With
  the same seed, all ``random'' draws are repeated identically, so
  results can be reproduced. Without it every knit gives different
  numbers.
\item
  \texttt{rpois(n,\ lambda)}: draws \texttt{n} random numbers from a
  Poisson distribution with mean \texttt{lambda}. The Poisson
  distribution describes counts of independent events, which is a
  natural model for the number of offspring.
\item
  \texttt{max(0,\ x)}: returns the larger of the two values, used here
  as a floor at 0.
\item
  \texttt{rbinom(n,\ size,\ prob)}: draws \texttt{n} random numbers from
  a binomial distribution, i.e.~the number of ``successes'' among
  \texttt{size} independent trials, each with success probability
  \texttt{prob}. Here: how many of the individuals are assigned to
  habitat B.
\item
  \texttt{rbind(x,\ y)}: ``row-bind'', appends \texttt{y} as a new row
  below \texttt{x}. Starting from \texttt{NULL}, the first
  \texttt{rbind} creates the matrix. This is how the time series grows
  by one row per generation.
\item
  \texttt{rownames()} / \texttt{colnames()}: set the labels of the rows
  / columns of a matrix. \texttt{(0:t\_max){[}1:length(...){]}} creates
  the generation numbers and cuts them to the actual number of rows.
\item
  \texttt{break}: leaves the \texttt{for} loop immediately (used for
  early stopping).
\item
  \texttt{\textbar{}}: element-wise logical ``or'' (here equivalent to
  \texttt{\textbar{}\textbar{}} since both sides are single values).
\item
  \texttt{as.data.frame()}: converts a matrix into a data frame, so
  columns can be accessed with \texttt{\$} (e.g.~\texttt{output\$B\$A}).
  \texttt{as.data.frame(NULL)} gives an empty data frame with 0 rows,
  which is how we detect runs that ended before the split.
\item
  \texttt{list(main\ =\ ...,\ B\ =\ ...,\ C\ =\ ...)}: returns several
  objects of different size in one named container.
\item
  \texttt{par()}: sets global graphics parameters. \texttt{mar} =
  margins (bottom, left, top, right) in lines;
  \texttt{mfrow\ =\ c(2,\ 1)} = a grid of 2 rows × 1 column of plots.
  \texttt{par()} returns the \emph{old} values, which we save in
  \texttt{old\_par}.
\item
  \texttt{on.exit(par(old\_par))}: registers code that runs
  automatically when the function ends, restoring the original graphics
  settings so that later plots are not affected.
\item
  \texttt{segments(x0,\ y0,\ x1,\ y1)}: draws a straight line between
  two points; used for the carrying-capacity lines that only span part
  of the x-axis.
\item
  \texttt{legend()}: adds a legend. \texttt{par("usr")} returns the plot
  region's coordinates (x-min, x-max, y-min, y-max), so
  \texttt{par("usr"){[}2{]}} and \texttt{{[}4{]}} are the top-right
  corner. \texttt{xpd\ =\ TRUE} allows drawing outside the plot region,
  \texttt{bty\ =\ "n"} removes the legend box,
  \texttt{xjust}/\texttt{yjust} set which corner of the legend is
  anchored at that point.
\item
  \texttt{seq\_along(x)} / \texttt{seq\_len(n)}: safe ways to create
  \texttt{1,\ 2,\ ...,\ length(x)} or \texttt{1,\ ...,\ n} (they return
  an empty sequence for empty input, unlike \texttt{1:length(x)}).
\item
  \texttt{if\ (...)\ a\ else\ b} used as an expression:
  \texttt{x\_max\ \textless{}-\ if\ (is.null(xlim\_max))\ ...\ else\ ...}
  assigns one of two values.
\item
  \texttt{plot.new()} + \texttt{title()}: create a blank plot and add a
  title, used as a placeholder when there is no data.
\item
  \texttt{invisible(NULL)}: returns \texttt{NULL} without printing it.
\item
  \texttt{paste0()}: joins strings without separator, used to put the
  starting numbers in the legend.
\item
  \texttt{repeat\ \{\ ...\ \}}: an endless loop that is left only with
  \texttt{break}. Used here because we do not know in advance how many
  attempts are needed to get enough valid runs.
\item
  \texttt{next}: skips the rest of the current loop iteration and goes
  to the next one.
\item
  \texttt{any(x)}: \texttt{TRUE} if at least one element of the logical
  vector \texttt{x} is \texttt{TRUE}.
\item
  \texttt{which(x){[}1{]}}: \texttt{which()} returns the positions of
  all \texttt{TRUE} elements; \texttt{{[}1{]}} takes the first, i.e.~the
  \textbf{first generation} at which the threshold was reached. Note
  this counts rows after the split (row 1 = first generation after the
  split).
\item
  \texttt{NA\_integer\_}: a missing value of integer type (a plain
  \texttt{NA} could create type conflicts when binding rows).
\item
  \texttt{results{[}{[}length(results)\ +\ 1{]}{]}\ \textless{}-\ ...}:
  appends an element to the end of a list.
\item
  \texttt{do.call(rbind,\ results)}: calls \texttt{rbind} with the list
  elements as separate arguments, i.e.~stacks all data frames in the
  list into one big table.
\item
  \texttt{\%\textgreater{}\%} (pipe): passes the result on the left as
  the first argument of the function on the right, so operations can be
  read top to bottom.
\item
  \texttt{group\_by(st,\ hp)}: tells \texttt{dplyr} to do the following
  calculations separately for each combination of \texttt{st} and
  \texttt{hp}.
\item
  \texttt{summarise(...)}: collapses each group into a single row.
  \texttt{n()} counts the rows in the group;
  \texttt{across(cols,\ mean)} applies \texttt{mean} to several columns
  at once; \texttt{.groups\ =\ "drop"} removes the grouping afterwards.
\item
  \texttt{pivot\_longer()}: turns several columns into two
  (\texttt{names\_to} gets the old column names, \texttt{values\_to} the
  values). One row per scenario and population results, as
  \texttt{ggplot2} needs..
\item
  \texttt{ggplot(data,\ aes(...))\ +\ geom\_*()}: \texttt{ggplot2}
  builds a plot in layers. \texttt{aes()} maps data columns to visual
  properties (x, y, colour, fill), and each \texttt{geom\_} adds a
  layer: \texttt{geom\_line}, \texttt{geom\_point},
  \texttt{geom\_errorbar} (error bars), \texttt{geom\_tile} (coloured
  rectangles), \texttt{geom\_text} (printed numbers).
\item
  \texttt{pmax()} / \texttt{pmin()}: \emph{parallel} maximum / minimum,
  applied element by element to a vector (unlike \texttt{max()}, which
  returns a single number). Used to keep error bars inside {[}0, 1{]}.
\item
  \texttt{scale\_y\_continuous(limits,\ breaks)}: sets the axis range
  and tick positions. \texttt{seq(from,\ to,\ by)} generates the tick
  values.
\item
  \texttt{labs()}: sets titles and axis/legend labels.
  \texttt{theme\_minimal()}: applies a clean plot theme.
\item
  \texttt{facet\_wrap(\textasciitilde{}\ population)}: splits the plot
  into one panel per value of \texttt{population}.
\item
  \texttt{scale\_fill\_viridis\_c(option\ =\ "C")}: a continuous colour
  scale (viridis palette), readable also for colour-blind people.
\item
  \texttt{factor(x,\ levels,\ labels)}: converts values into a
  categorical variable with a fixed order and custom display labels.
  \texttt{factor(st)} makes the split time a category so that tiles are
  evenly spaced.
\item
  \texttt{sprintf("\%.2f",\ x)}: formats a number as text with 2
  decimals.
\item
  \texttt{sort(unique(x))}: the distinct values of \texttt{x} in
  ascending order. code chunks for maybe later use:
\end{itemize}

\end{document}
